\def\thoigian{90}%--Thời gian
\de{ĐỀ THI GIỮA HỌC KỲ II NĂM HỌC 2023-2024}{THPT Nguyễn Tất Thành - TP HCM}

%Câu 1...........................
\begin{bt}%[1D6H1-2]%[Dự án đề kiểm tra Toán 10 GHKII NH23-24- Đoàn Minh Tâm]%[THPT Nguyễn Tất Thành - TPHCM]
Rút gọn biểu thức $P=\dfrac{(a b)^{2-\sqrt{5}} \cdot b^{\sqrt{5}-1}}{a^{3-\sqrt{5}}}$.
\loigiai{Ta có $P=\dfrac{(a b)^{2-\sqrt{5}} \cdot b^{\sqrt{5}-1}}{a^{3-\sqrt{5}}}=a^{2-\sqrt{5}} \cdot a^{-3+\sqrt{5}} \cdot b^{2-\sqrt{5}}\cdot b^{\sqrt{5}-1}=a^{-1} \cdot b=\dfrac{b}{a}$.}
\end{bt}

%Câu 2...........................
\begin{bt}%[1D6H2-1]%[Dự án đề kiểm tra Toán 10 GHKII NH23-24- Đoàn Minh Tâm]%[THPT Nguyễn Tất Thành - TPHCM]
	Biết $\log _a b=2$, $\log _c a=3$ ($a$, $b$, $c$ dương khác $1$). Tính $\log _{a b}(b c)$.
	\loigiai{
	Ta có $\log _{a b}(b c)=\dfrac{\log _a b c}{\log _a a b}=\dfrac{\log _a b+\log _a c}{\log _a a+\log _a b}=\dfrac{2+\dfrac{1}{3}}{1+2}=\dfrac{7}{9}$.
}
\end{bt}

%Câu 3...........................
\begin{bt}%[1D6N3-5]%[Dự án đề kiểm tra Toán 10 GHKII NH23-24- Đoàn Minh Tâm]%[THPT Nguyễn Tất Thành - TPHCM]
	Dân số thế giới sau $t$ năm kể từ năm $2020$ được tính bởi công thức\\ $P(t)=A \cdot(1+0{,}0105)^{t}$ (tỉ người). Dân số thế giới năm $2024$ ước tính $8{,}128$ tỉ người, dân số thế giới năm $2030$ là bao nhiêu tỉ người? (làm tròn kết quả đến hàng phần ngàn)
	\loigiai{
		Do năm $2024$ dân số ước tính là $8{,}128$ nên ta có
		$$ 8{,}128=A\cdot(1+0{,}0105)^{4} \Rightarrow A \approx 7{,}795. $$
		Vậy dân số ước tính năm $2030$ là
		$$ P(10)=7{,}795\cdot (1+0{,}0105)^{10} \approx 8{,}654. $$
	}
\end{bt}

%Câu 4...........................
\begin{bt}%[1D6N3-2]%[Dự án đề kiểm tra Toán 10 GHKII NH23-24- Đoàn Minh Tâm]%[THPT Nguyễn Tất Thành - TPHCM]
	Tìm tập xác định của hàm số $y=\log _2\left(4-3 x-x^2\right)$.
	\loigiai{Điều kiện $4-3x-x^2>0 \Leftrightarrow -4<x<1$.\\
		 Vậy tập xác định là $\mathscr{D}=(-4;1)$.}
\end{bt}

%Câu 5...........................
\begin{bt}%[1D6H4-4]%[Dự án đề kiểm tra Toán 10 GHKII NH23-24- Đoàn Minh Tâm]%[THPT Nguyễn Tất Thành - TPHCM]
	Giải các phương trình sau
	\begin{enumerate}
		\item $27^{2 x+1}=9^{x-2}$.
		\item $\log _2(x+1)=6-\log _2(5 x+1)$.
	\end{enumerate}
	\loigiai{
		\begin{enumerate}
			\item $27^{2 x+1}=9^{x-2} \Leftrightarrow 3^{6x+3}=3^{2x-4}\Leftrightarrow 6x+3=2x-4 \Leftrightarrow x=-\dfrac{7}{4}$.
			\item Điều kiện $\heva{&x+1>0 \\&5x+1>0} \Leftrightarrow \heva{&x>-1\\&x>-\dfrac{1}{5}} \Leftrightarrow x>-\dfrac{1}{5}$.\\
				 Phương trình tương đương $$ \log _2(x+1)(5 x+1)=6 \Leftrightarrow 5 x^2+6 x-63=0 \Leftrightarrow \hoac{&x=3\text{ (nhận)}\\ &x=-\dfrac{21}{5}\text{ (loại)}.}$$
		\end{enumerate}
	}
\end{bt}

%Câu 6...........................
\begin{bt}%[1D6H4-4]%[Dự án đề kiểm tra Toán 10 GHKII NH23-24- Đoàn Minh Tâm]%[THPT Nguyễn Tất Thành - TPHCM]
	Giải các bất phương trình sau
	\begin{enumerate}
		\item $9^{2 x+1} \leq 3^{x+2}$.
		\item $\log _2(x+1)>2 \log _4(5 x+1)$.
	\end{enumerate}
	\loigiai{
		\begin{enumerate}
			\item $9^{2 x+1} \leq 3^{x+2} \Leftrightarrow 3^{4 x+2} \leq 3^{x+2} \Leftrightarrow 4 x+2 \leq x+2 \Leftrightarrow x \leq 0$.\\
			Vậy $S=(-\infty;0]$.
			\item 	Điều kiện $\heva{&x+1>0 \\&5x+1>0} \Leftrightarrow \heva{&x>-1\\&x>-\dfrac{1}{5}} \Leftrightarrow x>-\dfrac{1}{5}$.
			$$ \log _2(x+1)>2 \log _4(5 x+1)
			\Leftrightarrow \log _2(x+1)>\log _2(5 x+1) \Leftrightarrow x+1>5 x+1 \Leftrightarrow x<0.
			$$
			So với điều kiện $S=\left(-\dfrac{1}{5} ; 0\right)$.
		\end{enumerate}
	}
\end{bt}

%Câu 7...........................
\begin{bt}%[1H8H2-2]%[1H8H1-3]%[Dự án đề kiểm tra Toán 10 GHKII NH23-24- Đoàn Minh Tâm]%[THPT Nguyễn Tất Thành - TPHCM]
	Cho hình chóp $S . A B C D$ có đáy là hình vuông cạnh $a$, $S A \perp(A B C D)$, $S A=\sqrt{2} a$, $A K$ là đường cao của $\triangle S A D$.
	\begin{enumerate}
		\item Chứng minh $A K \perp(S C D)$.
		\item Tính góc giữa hai đường thẳng $B K$ và $C D$.
	\end{enumerate}
	\loigiai{
		\immini{
			\begin{enumerate}
				\item Ta có $\heva{&CD \perp AD\\&CD\perp SA}\Rightarrow CD \perp (SAD)$.\\
				Mà $AK \subset (SAD) \Rightarrow AK \perp CD$.\\
				Vậy ta có $\heva{&AK \perp SD\\&AK\perp CD}\Rightarrow AK \perp (SCD)$.
				\item $(B K, C D)=(B K, A B)$ (do $CD \parallel AB$).\\
				Do $\heva{&CD\perp (SAD)\\&CD \parallel AB} \Rightarrow AB \perp (SAD)$.\\
				Mà $AK \subset (SAD) \Rightarrow AB \perp AK$ nên $\triangle ABK$ vuông tại $A$.\\
				Ta có $AK=\dfrac{AS \cdot AD}{\sqrt{SA^2+AD^2}}=\dfrac{a\sqrt{6}}{3}$.\\
				Vậy $\tan ABK = \dfrac{AK}{AB}=\dfrac{\sqrt{6}}{3}$.\\
				Nên $(B K, C D)=(B K, A B)=\widehat{ABK}\approx 39^\circ 13'$.
			\end{enumerate}
		}{
			\begin{tikzpicture}[line join=round,line cap=round,line width=.6pt,font=\footnotesize,scale=1]
				\coordinate (B) at (0,0);
				\coordinate (A) at (1.5,.8);
				\coordinate (C) at (4,0);
				\coordinate (D) at ($(C)-(B)+(A)$);
				\coordinate (S) at ($(A)+(90:4)$);
				\coordinate (K) at ($(S)!(A)!(D)$);
				\draw (B)--(C)--(D)--(S)--cycle (S)--(C);
				\draw[dashed] (A)--(D) (S)--(A)--(B) (A)--(K)--(B);
				\draw ($ (A)!5pt!(D)$)--($(A)!2!($($(A)!5pt!(D)$)!.5!($(A)!5pt!(S)$)$)$)--($(A)!5pt!(S)$);
				\draw ($ (K)!5pt!(D)$)--($(K)!2!($($(K)!5pt!(D)$)!.5!($(K)!5pt!(A)$)$)$)--($(K)!5pt!(A)$);
				\foreach \x/\y in {A/180, B/180, C/0, D/0, S/90,K/60}{\fill (\x) circle(1pt) ($(\x)+(\y:0.3cm)$) node{$\x$};}
			\end{tikzpicture}

		}
	}
\end{bt}
